\section{Relative Work}
\begin{figure*}[ht!]
	% 		\vspace{-0.2cm}  %
	% 		\setlength{\abovecaptionskip}{-0.0cm}   %
	% 		\setlength{\belowcaptionskip}{-0.5cm}
	\centering
	\includegraphics[width=1\textwidth]{figures/deploymemta-visual.png}
	%\caption{\textcolor{red}{a): Equipment Layout; Second col: vehicles-side image and point cloud with 3d annotation; third col: infrastructure-side image and point cloud with 3d annotation.}}
	\caption{a) Acquisition system with infrastructure sensors. b) Acquisition system with vehicle sensors. c) Infrastructure-view image and point cloud with 3d annotation. Paired vehicle-view and infrastructure-view information complement each other in the perspective of view. d) Vehicle-view image and point cloud with 3d annotation.}
	\label{fig:equipment}
\end{figure*}
\subsection{Autonomous Driving Datasets}
In recent years, an increasing number of autonomous driving datasets have been released and greatly promoted the development of autonomous driving research. Datasets like SYNTHIA ~\cite{ros2016synthia} and Cityscapes ~\cite{cordts2016cityscapes} mainly focus on 2D annotations for images. KITTI~\cite{Geiger2012KITTI} and nuScenes~\cite{Caesar2020NuScenes} are multi-modality datasets providing camera images as well as LiDAR point clouds. Nevertheless, all datasets mentioned above only provide data from a single-vehicle view. V2X-SIM~\cite{li2021learning} is an attempt to generate a multi-vehicle view dataset, but the dataset was generated by a simulator rather than captured from real scenarios. Compared with those datasets, our DAIR-V2X dataset is the first large-scale, multi-modality, multi-view dataset captured from real scenarios for VICAD, and contains data captured from the Vehicle-Infrastructure Cooperative view.
Tab.~\ref{tab: dataset comparison} shows the comparison of our dataset with the others. 
In our DAIR-V2X, we also provide a Repo3D~\cite{ye2022rope3d} dataset composed of multi-source infrastructure images and 3D annotations, for those who are interested in Mono3D object detection and domain adaptation.

% Different from KITTI where only frontal view objects are annotated, nuScenes~\cite{Caesar2020NuScenes} provides a full 
% 360$^{\circ}$ vehicle view. 
 

%It demonstrates that DAIR-V2X is the first large-scale, multi-modalityity, multi-view dataset for Vehicle-Infrastructure Cooperative Autonomous Driving.

%It is easy to notice that most published autonomous driving datasets only provide data from single-vehicle side. The KITTI~\cite{Geiger2012KITTI} is the pioneer among them, which contains 15K manually labeled images and point clouds. The nuScenes~\cite{Caesar2020NuScenes} is the first dataset to have full 360$^{\circ}$ field of vehicle view. The Waymo Open ~\cite{sun2020scalability} offers much more diverse data compared with the history. Different from above, the V2X-Sim dataset~\cite{ZHAO2019Detection} is a multi-vehicle dataset generated by simulator. There are also some other datasets concerning infrastructure side rather than vehicle side. For instance, the highD~\cite{krajewski2018highd} records highway information from drone and offers high quality labels. Tab.~\ref{tab: dataset comparison} shows the comparison of our dataset with others. It demonstrates that DAIR-V2X is the first large-scale, multi-modalityity, multi-view dataset for Vehicle-Infrastructure Cooperative Autonomous Driving.

% 

%There are also some other datasets providing video clips with image annotations from vehicle side. For example, the CityFlow\cite{tang2019cityflow} dataset contains 3.25 hours of videos and more than 200k bounding boxes. The Citysapes 3D dataset\cite{gahlert2020cityscapes} is an extension of the Cityscapes dataset\cite{cordts2016cityscapes} with both videos and 3D bounding box for 23 types of vehicles.

%Only a very few datasets facilitate other perceptions besides single-vehicle side. The highD\cite{krajewski2018highd} and inD\cite{bock2020ind} datasets are drone datasets that offers recordings which 2D labels. The V2X-Sim dataset\cite{Bonjorn20185GV2X} is a multi-agent dataset generated by simulator. 

%Tab.~\ref{tab: dataset comparison} shows the comparison of properties of our dataset and others. It demonstrates that ours is the first large-scale, multi-modalityity, multi-view dataset for Vehicle-Infrastructure cooperative Autonomous Driving.

% Please add the following required packages to your document preamble:
% \usepackage{graphicx

\subsection{3D Detection}
3D object detection serves as the prerequisite for the success of autonomous driving. Many techniques have been introduced and can be roughly classified into three categories. a) Image-based 3D Detection refers to methods that detect 3D objects directly from 2D images. ImVoxelNet\cite{rukhovich2021imvoxelnet} is a good example to make predictions from images. 
b) Pointcloud-based 3D Detection stands for manners that make 3D object detection merely from point clouds. PointPillars~\cite{lang2019pointpillars}, SECOND~\cite{yan2018second}, and 3DSSD~\cite{Yang2020ssd} are such approaches that achieve convincing detection results from point clouds. 
c) Multimodality-based 3D Detection uses both images and point clouds to make predictions. Pointpainting~\cite{Vora2020Point} and MVXNet~\cite{Sindagi2019MVX} are practices of fusing image and LiDAR features to predict 3D bounding boxes.
While 3D object detection has made great progress recently, there are still some tough problems that remain to solve such as blind spots and weak  long-distance perception. 
To explore how to utilize the infrastructure information to solve the problems mentioned above, 
we conduct VIC3D object detection based on our dataset proposed in this paper.

\subsection{Multi-Sensor Fusion}

% why multi-sensor
Multi-sensor fusion\cite{Wang2020Multi} is the integration of heterogeneous information collected by different sensors to alleviate the uncertainty and vulnerability of systems that rely on a single sensor.
%and improve perception results of autonomous driving vehicles. %Conventional sensors including radar, camera, LiDAR, ultrasonic, GPS and V2X possess different properties and capabilities. The selection of multi-sensor combination includes camera-LiDAR fusion  \cite{Chen2017Multi, Sindagi2019MVX, Wang2019Frustum, Wang2020Multi}, Radar-Camera fusion\cite{Chadwick2019Distant} and LiDAR-Radar fusion\cite{Wang2019radar}. 
Based on the fusion stage, multi-sensor fusion can be categorized into early fusion, intermediate fusion, and late fusion. a) In early fusion, raw data from different sensors are directly transferred and fused~\cite{gao2018early}. b) In intermediate fusion, intermediate representations like features extracted from the models are fused~\cite{Chen2017Multi, Sindagi2019MVX}. c) In late fusion, the prediction outputs like 3D information of the objects are fused~\cite{hag2018late}. 
VIC3D can be considered as a variant of the multi-sensor problem, so previous fusion methods can be taken into consideration to integrate the infrastructure information.
However, in addition to the multi-sensor fusion challenges, VIC3D faces difficulties caused by the temporal asynchrony problem and the data transmission constraint.

\subsection{V2X Cooperative Perception}
V2X aims to build a communication system between vehicles and other devices in a complex traffic environment. Current V2X research mainly focuses on V2V (Vehicle-to-Vehicle) and V2I (Vehicle-to-Infrastructure) area. V2VNet\cite{Wang2020V2VNetVC} is a pioneering work in V2V that broadcasts compressed intermediate features and propagates message received from nearby vehicles to generate motion forecasts.  Works of V2I~\cite{Cui2019Automatic, ZHAO2019Detection} leverage infrastructure LiDAR data to generate and broadcast detection results. 
However, none of these approaches have been verified on a dataset captured from real scenarios. This may cause a huge gap between theory and practice. Therefore, we release the DAIR-V2X dataset to boost further study in this field.